\section{Conclusions, limitations and plans}
\label{sec:conclusions}

\NHFOAM\ follows NHWAVE's pressure arrangement on OpenFOAM's unstructured,
moving $\sigma$-meshes. The tests of \citet{ma2012} were run with both models
on the same set-ups, from the simplest upwards, and they changed the code:
\begin{itemize}[nosep]
  \item \textbf{Time accuracy.} With one outer corrector the convection term
        was linearised with the flux of the previous time level, which made the
        scheme first order in time for nonlinear waves. The convecting flux is
        now extrapolated from the last two levels. A solitary wave of
        $H/h=0.5$ lost 6.4\% of its height in 150 m before and 1.0\% after the
        fix (NHWAVE 2.3\%), at no extra cost.
  \item \textbf{Convection.} The TVD scheme damps steep short waves even when
        they do not break; filtered central differencing
        (\code{filteredLinear2V 0.2 0}) is now recommended for non-breaking
        waves. It reduced the standing-wave error to 0.16--0.25\% for
        $kD\ge 2\pi$ (NHWAVE 0.6--4.8\%) and stayed stable in 40--65 s runs over
        the submerged bar. Breaking still needs the TVD scheme.
  \item \textbf{Consistency with NHWAVE.} After the fix, the phase of the waves
        over and behind the submerged bar agrees with NHWAVE's. The closer
        match of the measured phase before the fix came from the first-order
        lag, which slowed the waves by about 0.5\%.
  \item \textbf{Resolution.} On the elliptic shoal, with twice the paper's
        grid spacing and three layers, \NHFOAM's errors at the first sections
        (0.06 and 0.13) are close to NHWAVE's on the paper's grid (0.07 and
        0.11), while NHWAVE on the coarse grid loses 30\% of the incident wave
        height before the shoal.
  \item \textbf{Robustness.} A domain in which waves break anywhere needs the
        TVD scheme and the breaking model, even outside the region of
        interest: the shoal run with the filtered scheme failed when the waves
        broke in a shallow corner, and ran on with \code{limitedLinearV} and
        the hydrostatic front approximation.
  \item \textbf{Wetting and drying.} NHWAVE's mask and thin film carry over,
        but the implicit free surface needs a limiter on the outflows (which
        keeps the volume conserved, unlike NHWAVE's clip), and the collocated
        mesh needs the upwind (TVD) depth at the side faces and the upwind
        predicted flux in supercritical flow. With the centred depth of the
        mesh faces the fronts were too slow ($R/d = 0.41$ for the solitary
        wave); with both changes the run-up is 0.531 (NHWAVE 0.512, lab
        0.49--0.55), still water stays at rest on a beach, and the
        parabolic basin is reproduced to 0.9 mm rms. A dam break on a dry
        bed showed that both models treat the film as a wet bed and follow
        Stoker's solution for it, not Ritter's. Parallel runs give the same
        run-up with a processor boundary in the swash zone (0.540 against
        0.531).
  \item \textbf{Breaking.} The HFA alone flags only the fronts of the
        bores. A persistence threshold that falls with the age of the
        breaking event and the breaking eddy viscosity of
        \citet{kennedy2000}, with horizontal mixing, nearly halve the
        surf-zone errors of the wave height on both spilling-breaker beaches (mean
        error $+6.1 \to +3.5$ cm on the beach of \citet{ting1994}, NHWAVE
        $+4.8$; rms $2.98 \to 1.54$ cm on that of \citet{hansen1979}), raise
        the set-up, and keep the measured shape of the undertow, which is
        30\% too strong (NHWAVE 60\%). Vertical mixing through the layer
        faces flattens the undertow and dissipates less. In 3D, columns
        must join a breaker at the threshold of a new breaker, or the
        breaking region creeps along oblique crests into deep water.
  \item \textbf{Longshore current.} Test 4.7 of \citet{ma2012}, in 3D
        with periodic lateral boundaries, wetting and drying, log-law
        friction and the Smagorinsky model, runs with both models on the
        same coarse grid. Both put the current at about the measured
        position (\NHFOAM: 0.265 m/s at 1.08 m from the shoreline, NHWAVE:
        0.250 m/s at 1.26 m, lab 0.39 m/s at 1.21 m) and give two thirds of
        its strength. With the HFA alone the current is 0.19 m/s and too
        close to the shoreline, as \citet{ma2012} found without turbulent
        diffusion.
  \item \textbf{Breaking parameters and closures.} Varied one at a time,
        the parameters of the default model show a shallow optimum near
        their published values: no change of $\alpha$, $\beta_F$ or
        $\delta_b$ moves the surf-zone errors by more than 8\%, and none
        removes the slow loss of height just after the breaking point.
        The kinematic onset with the crest speed estimated from the front
        face and $B_k = 0.70$ breaks at the measured place on the beach of
        \citet{ting1994}, with the smallest surf-zone error so far ($+2.4$
        cm), and 0.7 m early on that of \citet{hansen1979}. The stabilised
        $k$--$\omega$ model of \citet{larsen2018} keeps the turbulence of
        the nearly potential flow 30--60\% below that of the unstabilised
        model, but it does not replace the breaking model: without it the
        waves break 1.8 m late, and with the HFA it adds little
        dissipation.
  \item \textbf{Wave input.} The boundary wavemaker makes the target
        waves to within 0.5\%, and its active absorption keeps a closed
        flume steady (incident amplitude $1.03\,a$, against a beat between
        $0.1\,a$ and $1.1\,a$ without it); as an absorbing boundary it
        reflects 0.3--1.7\% of waves with $kh \le 1$ when given the wave's
        phase speed. Irregular JONSWAP waves keep the target spectrum,
        with $H_{m0}$ within $-4\%$ and $+2\%$ over 8 peak wavelengths.
  \item \textbf{Long runs.} The drift of the phase over the submerged bar
        (0.21 s in 20 s) came from the missing return current of the
        wave-making zone. With the return current, now the default, or the
        boundary wavemaker, the runs are steady, with an error of 1.0 cm
        in every window, and \NHFOAM's waves arrive 0.12 s early behind the
        bar crest, like those of NHWAVE and REEF3D::NHFLOW.
  \item \textbf{Unstructured meshes.} On triangulated strips the breaking
        solitary wave runs up to $R/d = 0.49$ (structured 0.531, lab
        0.49--0.55) and its backwash is slower; on a mesh refined around a
        vertical cylinder the diffracted waves are within 7\% of the theory
        of \citet{maccamy1954}.
  \item \textbf{Reproducibility.} Restarts continue the run to round-off,
        parallel runs reproduce serial ones with breakers across processor
        boundaries ($4\times10^{-7}$ m), and a regression suite of ten
        short cases checks both against stored references. The landslide
        rerun with its beach and wetting and drying showed that the wall
        was not the cause of its large trailing waves.
\end{itemize}

\paragraph{Limitations.}
\begin{itemize}[nosep]
  \item \textbf{Wetting and drying} has been tested in 2DV and on
        triangulated strips. The film acts as a wet bed, so thin swash flows
        run slower than over a truly dry bed (in NHWAVE as well), and
        puddles of a few millimetres above the film stay on the beach in
        the backwash; on triangles the run-up is 8\% lower than on the
        structured mesh and the backwash slower. A narrow cell at the
        shoreline raises the errors. It needs the interface pressure
        layout.
  \item \textbf{Breaking} has been tested on two spilling-breaker beaches
        and the longshore current. With the eddy viscosity the decay inside
        the surf zone is close to the measured one, but just after the
        breaking point the models lose too little height (3.5--4 cm too
        much at $x = 8$ m on the beach of \citet{ting1994}), and no
        parameter of the model removes this. The kinematic onset with
        $B_k \approx 0.7$--0.8 is promising but has been tested on the two
        2DV beaches only. The crest-speed estimate does not look across
        processor boundaries.
  \item \textbf{Turbulence closure.} The stabilised $k$--$\omega$ model
        gives the undertow its measured strength but, with the HFA, half
        the measured turbulent kinetic energy and little dissipation; it
        does not replace the breaking eddy viscosity.
  \item \textbf{Wave generation} is linear (no second-order correction
        for the bound long waves of wave groups). The absorbing boundary
        is uniform over the depth, so waves with $kh > 1.5$ are reflected
        by 9--13\%. Directional spreading has not been tested in a basin.
        Irregular waves lose 3--5\% of $H_{m0}$ over 12 m in the peak band,
        for a reason not yet identified.
  \item \textbf{Landslide trailing waves} are 1.5--3 times the measured ones
        (Sections~\ref{sec:seabed} and \ref{sec:ma-landslide}); the beach is
        not the cause.
  \item \textbf{OpenFOAM version.} The code was compiled and run with
        v1912 only; its API calls were checked against the v2512 headers.
\end{itemize}

\paragraph{Plans.} In order of priority:
\begin{enumerate}[nosep]
  \item breaking: the transition just after the breaking point; the
        kinematic onset ($B_k \approx 0.7$--0.8) on more beaches and on the
        longshore current, also on the grid of \citet{ma2012};
  \item second-order wave generation, and a directional-spreading test in
        a basin;
  \item the trailing waves of the landslide (slide motion, dissipation);
  \item a build and the regression suite with a current OpenFOAM release;
  \item depth-integrated granular slides for landslide tsunamis;
  \item coupling to \code{interFoam} and \code{sedInterFoam} through relaxation
        zones, and hybrid meshes around structures.
\end{enumerate}
